\chapter*{Samenvatting}
\setheader{Preface}

In samenwerking met loods M en haar stakeholders, zal een ontwerp gerealiseerd worden voor het historische transportschip Jansje. Hierbij is de harde eis dat de aandrijving in het nieuwe ontwerp draait op volledig hernieuwbare energie. Om tot dit ontwerp te komen, is eerst gekeken naar het vaarprofiel, en welke technische parameters hieruit voortvloeien. Door middel van een interview met de waarnemend eigenaar van Jansje is het operationeel profiel vastgesteld. 
\\\\
Daarna is er onderzoek gedaan naar de verschillende mogelijkheden voor energieopslagsystemen voor hernieuwbare energie. Dit is gedaan door middel van een uitgebreide literatuurstudie, een model, technische criteria en maatschappelijk relevante criteria. Er is gekeken naar zowel bestaande als nog experimentele techniek. De uitkomst hiervan is dat een combinatie van experimentele zeezoutbatterijen met lithium-ion batterijen de meest geschikte energiebron zal zijn. 
\\\\
Als laatst is onderzocht welke aandrijving het beste past bij het operationeel profiel van Jansje en de gekoze energiebron. Hieruit is voortgekomen dat een elektrische motor het meest geschikt is. Omdat er vele elektromotoren beschikbaar zijn, is ervoor gekozen om niet één elektromotor aan te bevelen, maar in plaats daarvan een set criteria mee te geven, waarmee er ten tijde van de daadwerkelijke implementatie de keuze gemaakt kan worden.



